\chapter{经典积分定理}
\label{chap:classical-integral-theorems-solutions}

\begin{solution}{green-simple-domain-proof}
设 $D=\{(x,y):a\le x\le b,\ \alpha(x)\le y\le\beta(x)\}$。沿正向边界积分
$P\,dx$ 时，两条竖边没有贡献，上、下图线方向相反，故
\[
 \int_{\partial D}P\,dx
 =\int_a^b[P(x,\alpha(x))-P(x,\beta(x))]dx
 =-\iint_D P_y\,dxdy.
\]
若又写成
$D=\{(x,y):c\le y\le d,\ \gamma(y)\le x\le\delta(y)\}$，同理
\[
 \int_{\partial D}Q\,dy
 =\int_c^d[Q(\delta(y),y)-Q(\gamma(y),y)]dy
 =\iint_D Q_x\,dxdy.
\]
两式相加得到 Green 公式。端点重合处只影响有限个点，不改变积分。
\end{solution}

\begin{solution}{green-annulus-orientation}
圆环的正向边界由外圆逆时针和内圆顺时针组成，使区域始终位于行进方向左侧。
例如取 $P=-y/2,Q=x/2$，则 $Q_x-P_y=1$。半径为 $R>r$ 的圆环上
\[
 \int_{\partial D}\frac{-y\,dx+x\,dy}{2}
 =\pi R^2-\pi r^2=|D|.
\]
若把内圆也取逆时针，左端会变成两圆面积之和，说明内边界定向不能省略。
\end{solution}

\begin{solution}{green-area-formula}
在 Green 公式中依次取 $(P,Q)=(0,x)$、$(-y,0)$ 和
$(-y/2,x/2)$，均有 $Q_x-P_y=1$。因此
\[
 |D|=\int_{\partial D}x\,dy
 =-\int_{\partial D}y\,dx
 =\frac12\int_{\partial D}(x\,dy-y\,dx).
\]
这些公式均要求边界按正向取向。
\end{solution}

\begin{solution}{green-internal-boundary-cancellation}
把区域分成有限片后，每条内部公共弧分别作为相邻两片边界出现。两片的正向
规定使这两次走向相反，因而同一连续一次形式在线积分中逐项抵消。若三片在
一个点相交，该点对一维 Riemann 积分没有贡献；公共弧仍成对归属。相加后
只留下整体外边界，右端的面积积分则由有限可加性合并。
\end{solution}

\begin{solution}{gauss-box-direct}
对盒 $[a,b]\times[c,d]\times[e,f]$，$x=b$ 与 $x=a$ 两面的外向通量之和为
\[
 \int_c^d\int_e^f(F_1(b,y,z)-F_1(a,y,z))\,dzdy
 =\iiint_Q\partial_xF_1\,dV.
\]
$y$、$z$ 两对面同理给出 $\iiint_Q\partial_yF_2$ 与
$\iiint_Q\partial_zF_3$。六面相加即为 $\iiint_Q\operatorname{div}F$。每一步只用一元
微积分基本定理和 Fubini 定理。
\end{solution}

\begin{solution}{gauss-volume-formula}
令 $F(x)=x/3$，则 $\operatorname{div}F=1$，所以
\[
 |\Omega|=\frac13\int_{\partial\Omega}x\cdot n\,dS.
\]
半径 $R$ 的球上 $x\cdot n=R$，得 $|B_R|=4\pi R^3/3$。对底半径 $R$、高
$h$ 的直圆锥，把原点取在锥顶，侧面母线方向与 $x$ 相切，故
$x\cdot n=0$；只有底面贡献，得到体积 $\pi R^2h/3$。
\end{solution}

\begin{solution}{gauss-inverse-square-flux}
在球面 $S_R$ 上，$n=x/R$，故
$F\cdot n=(x/\norm{x}^3)\cdot(x/R)=1/R^2$。通量为
$4\pi R^2/R^2=4\pi$。虽然 $\operatorname{div}F=0$ 在 $x\ne0$ 处成立，包含原点的球
内 $F$ 不定义，更不属于 $C^1$，不能直接应用散度定理。对球壳应用定理，
内外球面通量相消，恰好说明通量与半径无关。
\end{solution}

\begin{solution}{stokes-parameter-expansion}
写 $F=(P,Q,R)$、$X=(x,y,z)$。链式法则展开
\[
 \partial_u(F(X)\cdot X_v)-\partial_v(F(X)\cdot X_u).
\]
含 $X_{uv}$ 的两项相消；余项按 $P,Q,R$ 的偏导归组，得到
\[
 (R_y-Q_z)(y_uz_v-z_uy_v)
 +(P_z-R_x)(z_ux_v-x_uz_v)
 +(Q_x-P_y)(x_uy_v-y_ux_v),
\]
这正是 $\operatorname{curl}F(X)\cdot(X_u\times X_v)$。
\end{solution}

\begin{solution}{stokes-hemisphere-disk}
只要两个定向曲面具有同一条正向边界，Stokes 公式都给出
$\int_C F\cdot dr=\int_S\operatorname{curl}F\cdot n\,dS$。例如取
$F=(-y/2,x/2,0)$，则 $\operatorname{curl}F=(0,0,1)$。单位圆盘上通量为 $\pi$；上半球
上为 $\int_S n_3\,dS$，这是其向 $xy$ 平面的投影面积，也为 $\pi$。两种
曲面得到同一环流，所依赖的是公共边界及相容定向。
\end{solution}

\begin{solution}{curl-gradient-divcurl}
若 $\varphi\in C^2$，则
\[
 \operatorname{curl}\nabla\varphi=(\varphi_{zy}-\varphi_{yz},
 \varphi_{xz}-\varphi_{zx},\varphi_{yx}-\varphi_{xy})=0.
\]
若 $F=(P,Q,R)\in C^2$，展开 $\operatorname{div}\operatorname{curl}F$ 后得到六个二阶混合偏导，
每一对由 Schwarz 定理相消。$C^2$ 条件保证相关混合偏导连续并可交换。
\end{solution}

\begin{solution}{path-independent-potential-construction}
固定 $x_0$，定义 $\varphi(x)$ 为从 $x_0$ 到 $x$ 的任一路径积分；路径无关使
定义良好。对坐标方向 $e_i$ 和充分小 $h$，可把到 $x+he_i$ 的路径取为先到
$x$ 再走线段，故
\[
 \frac{\varphi(x+he_i)-\varphi(x)}h
 =\frac1h\int_0^hF_i(x+te_i)dt\longrightarrow F_i(x).
\]
因此所有偏导为 $F_i$；$F$ 连续又使这些偏导连续，故
$\varphi\in C^1$ 且 $\nabla\varphi=F$。
\end{solution}

\begin{solution}{punctured-plane-circulation}
令 $F(x,y)=(-y/(x^2+y^2),x/(x^2+y^2))$。在穿孔平面直接求导得
$Q_x-P_y=0$。单位圆参数 $(\cos t,\sin t)$ 上
$F=(-\sin t,\cos t)$，与切向量相同，所以
\[
 \int_C F\cdot dr=\int_0^{2\pi}1\,dt=2\pi.
\]
因此 $F$ 不可能是穿孔平面上的全局梯度场；局部无旋并不足以排除绕孔环流。
\end{solution}

\begin{solution}{star-shaped-curl-free-potential}
不妨星心为零。令
$\varphi(x)=\int_0^1F(tx)\cdot x\,dt$。对第 $j$ 个坐标求导，
\[
 \partial_j\varphi(x)=\int_0^1\left[F_j(tx)+t\sum_i x_i\partial_jF_i(tx)\right]dt.
\]
无旋给出 $\partial_jF_i=\partial_iF_j$，故被积式为
$\frac d{dt}[tF_j(tx)]$。积分得到 $F_j(x)$。所以
$\nabla\varphi=F$；星形性保证整条线段 $tx$ 留在定义域中。
\end{solution}

\begin{solution}{classical-theorem-hypothesis-audit}
本章 Green 公式要求平面区域可有限分割为简单区域，边界分段 $C^1$，场在
邻域中 $C^1$；Gauss 公式要求三维区域可作坐标简单分割，边界分片光滑并
定向；Stokes 公式要求定向分片正则曲面及相容的分段光滑边界。连续但无
连续偏导的场、分形边界、不可定向曲面以及含场奇点的区域均不由这些证明
覆盖。某些结论可在更弱条件下成立，但需要测度论、弱导数或几何测度论工具，
不能从本章的 Riemann 分割论证直接推出。
\end{solution}
